\section{Use failures to locate the changed operation}
Change the low bit of the first check-message byte from hexadecimal \texttt{31} to
\texttt{30}. The visible string changes from \texttt{123456789} to
\texttt{023456789}; the finalized checksum changes from \texttt{cbf43926} to
\texttt{dc8f2d65}. Trace both states through the first eight bit steps. The first
state difference arises before any later bytes are read, so the investigation need
not begin with file handling or table lookup.

Then restore the byte and change one table entry that the message actually indexes.
Compare the bit and table methods. A disagreement between the two methods can locate
an error in the table path, but agreement alone is insufficient: both routines could
share the same wrong polynomial or boundary convention. Keep independently specified
check values and a model with a different computational orientation.

\subsection*{Exercise: Plan the verification domain}
List cases that test arithmetic width, every possible single input octet, all table
entries, an unaligned source, an empty update, offsets within backing arenas, and
addresses above 32 bits. Add repeated calls that alternate methods and supplied return
continuations. Explain which properties final-memory comparison cannot establish when
a routine writes a byte and restores it before returning.

\subsection*{Exercise: Distinguish detection from trust}
Two different messages can have the same 32-bit result because there are more possible
messages than possible checksums. Explain why finding accidental corruption and
establishing who authored a file are different requirements. A stronger unkeyed hash
in the next project still needs a trusted expected digest or an authentication scheme
when the sender's identity matters.
